\par\noindent\textbf{(1)} 设 \(\Phi_{\rm in}=-E_{\rm in}r\cos\theta\)，\(\Phi_{\rm out}=-E_0r\cos\theta+C\cos\theta/r^2\)，前者原点正则、后者远场恢复外场。电势连续与 \(\varepsilon\partial_r\Phi\) 连续给两条独立线性边界方程。
\par\noindent\textbf{(2)} 由连续性求 \(E_{\rm in}=3\varepsilon_2E_0/(\varepsilon_1+2\varepsilon_2)\)，\(C=a^3E_0(\varepsilon_1-\varepsilon_2)/(\varepsilon_1+2\varepsilon_2)\)。将偶极电势 \(p\cos\theta/(4\pi\varepsilon_2r^2)\) 与 \(C\) 项比较，得 \(p=4\pi\varepsilon_2a^3(\varepsilon_1-\varepsilon_2)E_0/(\varepsilon_1+2\varepsilon_2)\)。
\par\noindent\textbf{(3)} 若球内的绝对极化定义为 \(\boldsymbol P=(\varepsilon_1-\varepsilon_0)\boldsymbol E_{\rm in}\)，则球表面束缚电荷密度 \(\sigma_b=\boldsymbol P\cdot\hat r=(\varepsilon_1-\varepsilon_0)E_{\rm in}\cos\theta\)。注意包围介质本身有极化时，两侧束缚面电荷应分别计算，而净界面电荷依两侧 \(\boldsymbol P\) 跳跃。
\par\noindent\textbf{(4)} 若 \(\varepsilon_1\to\infty\)，则 \(E_{\rm in}\to0\)、\(p\to4\pi\varepsilon_2a^3E_0\)，为完美导体球结果；若 \(\varepsilon_1=\varepsilon_2\)，则 \(C=p=0\)、\(E_{\rm in}=E_0\)，无介电界面散射。若 \(\varepsilon_1<\varepsilon_2\)，诱导偶极与远场方向相反。
